\section{让助手接上你的思路}
\subsection{先接好模型}
打开\menu{设置 → 个人 AI 配置}，新建一个命名连接，填写接口地址、API Key 和模型名称。一个连接可以填写多个模型，每行一个；不同服务商可以分别建连接。接口格式能识别时自动选择，不能识别时，按服务商文档选择 OpenAI 或 Anthropic 等对应格式。

提问框里的模型菜单可以切换本轮使用的模型。密钥属于你自己的账号，不与其他成员共享。各模型的上下文长度也可以单独填写：窗口越大，能携带的上下文越多，同时可能增加单次请求用量。若接口提示上下文过长，对话中的卡片可以直接调整；用量图在个人 AI 配置里按当天、七天或一个月查看。

\subsection{讨论与修改都可以}
直接在助手里说你正在做什么，例如：\emph{我在写实验设计，想让这一段更清楚地说明对照条件。}不必先想好一条完美提示词。只是问问题时，助手可以给文字解释，不一定产生修改。

选中文字后会出现一个小输入框。选区是定位线索，助手还可以读取当前文件、按需查找项目里的其他文件与图片，并提出跨文件建议。把关键要求说清楚，比反复粘贴整篇内容更有效。

\tryit{选中上一节的示例句，试着问：“保留意思，把这句话写得具体一点；如果信息不够，先问我。”}

建议会显示红色删除与绿色新增。你可以先改候选文字，再逐段\menu{接受}或\menu{拒绝}；需要时也可以全部接受。接受前正文不会改变，已经接受的修改可以撤销。若你在等待时又改过同一段，先核对冲突提示，再决定采用哪份文字。

\subsection{从一句话走到整个项目}
助手能按需列目录、搜索内容、读文件，提出新建文件或目录、移动文件、修改编译设置等项目操作。比如：\emph{把引言独立为一个章节，并补上主文件中的引用。}先让它看真实文件，再检查提出的文字与文件操作，接受后才会写入。

图片需要模型具备视觉能力。如果当前模型不支持，系统会尝试去掉图片继续；要讨论图中的细节，请切换到能看图的模型。

\menu{设置 → Agent}管理搜索连接、技能和项目记忆。搜索的接口类型、endpoint 与 key 按搜索服务商填写；助手会在需要时调用。技能可以整理成反复使用的写作流程，也可以从网页导入。技能里的文字不会因此获得额外的系统权限。项目记忆适合放研究目标、术语约定与长期偏好，不必把整篇论文重复写进去。

\tryit{试着说：“先梳理本项目的文件结构，再告诉我哪些章节适合补一个例子；暂时只讨论。”这样可以练习让助手先了解作品，再进入修改。}
